% Editable CV source. Compile with pdfLaTeX (Overleaf compatible).
% Content matches the research CV. Layout is recreated in LaTeX.
% Edit the text below; adjust margins and spacing here if content grows.
\documentclass[11pt,letterpaper]{article}
\usepackage[utf8]{inputenc}
\usepackage[T1]{fontenc}
\usepackage[margin=0.65in,top=0.45in,bottom=0.5in]{geometry}
\usepackage{xcolor,enumitem,microtype,fancyhdr}
\usepackage{lmodern}
\usepackage[colorlinks=true,urlcolor=roleblue]{hyperref}
\definecolor{sectionred}{HTML}{900028}
\definecolor{roleblue}{HTML}{1749A5}
% PORTFOLIO: paste your real GitHub Pages URL between the braces below.
% Example format only: https://YOUR-USERNAME.github.io/portfolio/
% Leave empty until your website is published. Empty means no broken hyperlink.
\newcommand{\portfolioURL}{https://mojtabahouballah.github.io}

\newcommand{\emptyPortfolio}{}

\newcommand{\portfoliolink}{\ifx\portfolioURL\emptyPortfolio Portfolio: link to add\else\href{\portfolioURL}{Portfolio website}\fi}

\pagestyle{fancy}

\fancyhf{}
\renewcommand{\headrulewidth}{0pt}
\fancyfoot[R]{\footnotesize Mojtaba Houballah\quad\thepage}
\setlength{\parindent}{0pt}
\setlength{\parskip}{3pt}
\setlength{\emergencystretch}{2em}
\setlist[itemize]{leftmargin=14pt,label=\textbullet,itemsep=2pt,topsep=3pt,parsep=0pt,partopsep=0pt}
\newcommand{\cvsection}[1]{\par\addvspace{10pt}{\color{sectionred}\bfseries\fontsize{11pt}{13pt}\selectfont\MakeUppercase{#1}}\par\nobreak\vspace{3pt}\hrule height 0.4pt\relax\vspace{5pt}}
\newcommand{\cvrole}[1]{\par\addvspace{5pt}{\color{roleblue}\bfseries #1}\par\nobreak}
\hypersetup{pdftitle={Mojtaba Houballah CV},pdfauthor={Mojtaba Houballah}}
\begin{document}
\fontsize{11pt}{13.5pt}\selectfont
\begin{center}
{\fontsize{25pt}{29pt}\selectfont\MakeUppercase{Mojtaba Houballah}}\\[5pt]
Numerical Modeling Scientist\quad |\quad PhD in Informatics\\[3pt]
Paris region, France\quad |\quad \href{mailto:mojtaba.houballah@gmail.com}{mojtaba.houballah@gmail.com}\\[3pt]
\portfoliolink
\end{center}
% Profile
\cvsection{Profile}
Numerical modeling scientist with a PhD in Informatics and a background in optimization. Develops, calibrates, and evaluates process-based environmental models driven by climate and meteorological inputs. Experienced in scientific computing, large simulation workflows, site-based validation, and uncertainty analysis. Interested in applying these methods to renewable-energy resource and yield assessment.\par

% Technical expertise
\cvsection{Technical expertise}
\textbf{Numerical models: }Mathematical formulation, process-based simulation, implementation, and validation.\par
\textbf{Meteorological and climate data: }Forcing preparation, quality control, large datasets, and site observations.\par
\textbf{Calibration and uncertainty: }Optimization, sensitivity analysis, data assimilation, and error propagation.\par
\textbf{Scientific computing: }Python, Fortran, R, shell scripting, Linux, and high-performance computing.\par
\textbf{Collaboration and supervision: }Interdisciplinary research, technical presentations, and student mentoring.\par

% Professional experience
\cvsection{Professional experience}
\cvrole{Postdoctoral Researcher\hfill Sep 2023 - Aug 2026}
\textit{LSCE / IPSL, France}\par
\begin{itemize}
\item Developed and evaluated ORCHIDEE-CROP, integrating STICS formulations into a coupled land-surface model with carbon, nitrogen, water, and energy processes.
\item Developed Fortran, Python, R, and shell workflows for meteorological and climate-data preparation, quality control, model execution, and analysis of large numerical outputs on HPC systems.
\item Implemented a coupled nitrogen-water stress scheme; diagnosed simulation behavior across European sites with contrasting climate and management conditions.
\item Applied ORCHIDAS to parameter optimization, sensitivity analysis, and assessment of parameter transferability; evaluated simulations against observations.
\item Supervised a student project on maize parameter optimization with ORCHIDAS in 2026.
\end{itemize}
\cvrole{Postdoctoral Researcher\hfill Oct 2021 - Jul 2023}
\textit{ESE, Université Paris-Saclay, France}\par
\begin{itemize}
\item Developed and calibrated AMG-Forest, a process-based soil-carbon model evaluated against observations from the French RENECOFOR forest-monitoring network.
\item Analyzed parameter sensitivity and model behavior using soil-carbon observations and Rock-Eval-derived indicators.
\item Compared performance across sites and contributed to a framework for coupling with CASTANEA.
\end{itemize}

% Page two begins here. Remove this break if you prefer automatic pagination.
\newpage

% Professional experience
\cvsection{Professional experience}
\cvrole{Postdoctoral Researcher\hfill Jun 2020 - Jul 2021}
\textit{INRAE-SILVA, Nancy, France}\par
\begin{itemize}
\item Combined heterogeneous institutional datasets into regional resource-flow models using data reconciliation.
\item Assessed measurement uncertainty and error propagation in regional biomass and carbon estimates.
\item Compared resource-use scenarios with carbon-accounting tools in collaboration with researchers and regional stakeholders.
\end{itemize}
\cvrole{Doctoral Researcher\hfill Dec 2016 - Dec 2019}
\textit{Université Clermont Auvergne / INRAE, France}\par
\begin{itemize}
\item Developed dynamic models of environmental and management systems, linking physical processes, decisions, and economic constraints.
\item Applied optimization, control theory, and viability theory to identify management strategies that remain feasible under changing constraints.
\item Designed long-term simulations and multicriteria assessments, applied to forest management; produced three first-author peer-reviewed publications.
\end{itemize}

% Education
\cvsection{Education}
\textbf{PhD in Informatics, 2019: }Université Clermont Auvergne, France. Mathematical modeling of environmental systems using viability theory.\par
\textbf{Master II in Optimization, 2016: }Université Paris-Saclay, France.\par
\textbf{Master I in Mathematics, 2015: }Lebanese University, Beirut.\par
\textbf{BSc in Pure Mathematics, 2014: }Lebanese University, Nabatieh.\par

% Technical communication and mentorship
\cvsection{Technical communication and mentorship}
\begin{itemize}
\item Supervised Ernest Djmodi on crop parameter optimization using ORCHIDAS (2026) and Jean Weber on territorial wood-flow quantification (2021).
\item Presented ORCHIDEE-CROP development at the STICS Conference, Belgium (2026), and forest soil-carbon modeling at the EGU General Assembly, Austria (2023).
\end{itemize}

% Selected research outputs
\cvsection{Selected research outputs}
\begin{itemize}
\item Houballah, M. et al. Simulation of decadal variations in soil organic carbon stocks in French forests. \textit{Geoderma}, 2026. \href{https://www.sciencedirect.com/science/article/pii/S0016706126002545}{Published article}.
\item Houballah, M., Cordonnier, T., Mathias, J.-D. Maintaining or building roads? An adaptive management approach for preserving forest multifunctionality. Forest Ecology and Management, 537, 120957, 2023.
\item Houballah, M. et al. Development of ORCHIDEE-CROP v1.0: a nitrogen-explicit agro-land-surface model and its evaluation across European wheat sites. Manuscript in preparation.
\end{itemize}

% Modeling frameworks
\cvsection{Modeling frameworks}
ORCHIDEE-CROP, ORCHIDEE, STICS, ORCHIDAS, AMG-Forest, AMG, Century, CASTANEA, Carbon Accounting Tool.\par
\end{document}
